\section{Field behavior of the final quartic transformations}
\label{sec:final-quartic-fields-v2}

The last transformations improve the quartic tables without
new arithmetic input.  They give the preferred translated
1315-coordinate quartic, with coefficient height
$64\kappa$ for $\kappa=3^{2^{28}}$, and the preferred ordinary
integer-coefficient quartic with 1335 coordinates and height
less than $2^{1020}$.  Their field behavior follows precisely
from the type of transformation involved.  The four preceding
coefficient tables and their certificates remain unchanged.

\subsection{The affine translation works over every field of characteristic zero}

Put $S=3^{2^{27}}$, so $S^2=\kappa$.  The private anchor
coordinates $u=X_{55}$, $v=X_{56}$, $w=X_{57}$ in the
1315-coordinate table occur only in the row
\[
              -3A^2+16u^2+96v^2-32w^2+32S^2,
                    \qquad A=X_{52}+X_{53}-X_{54}.
\]
Make the rational affine substitution
\begin{equation}\label{eq:final-anchor-affine-map-v2}
 u=S+\frac{U-3V+3W}{4},\qquad
 v=\frac S2+\frac V4,\qquad
 w=\frac{3S}{2}+\frac W4.
\end{equation}
The new row is exactly
\[
 -3A^2+U^2-6UV+6UW+15V^2-18VW+7W^2+8SU.
\]
The inverse map is
\[
       U=4u+12v-12w+8S,\qquad V=4v-2S,
                                      \qquad W=4w-6S.
\]
Both maps are defined over $\mathbb Q$, with $S$ an ordinary
fixed integer.  Therefore the substitution is a bijection of
witness tuples over every characteristic-zero field.  It
preserves the literal polynomial's input projection, as well
as that of the residual system.  No use of an ordering is
required for this assertion.

For the generic $\mathbb Q(i)$ section of the preceding
1315-coordinate quartic, the old private anchor coordinates
are all zero.  Their inverse images are
\[
                   (X_{55},X_{56},X_{57})=(8S,-2S,-6S).
\]
Keep the Gaussian coordinates at indices 72 and 73, and set
all remaining witnesses to zero.  The restriction is the
same $R(T)+a^2+b^2$ as before, with $\kappa=S^2$.
Choosing $a=(1-R)/2$ and $b=i(1+R)/2$ proves that the
translated literal quartic also accepts every input over
every characteristic-zero field containing $i$.

\subsection{Eliminating a private squaring coordinate preserves the system}

In the preceding 1339-coordinate table, the coordinate
$z=X_{81}$ occurs only in the two rows
\[
                       X_{80}^2-z,
                         \qquad R_0+32z,
\]
where $R_0$ is independent of $z$.  Delete the first row and
$z$, and replace the second row by $R_0+32X_{80}^2$.
This gives 1338 coordinates and 1295 quadratic residuals.
Over every field, solutions of the new residual system lift
uniquely by $z=X_{80}^2$, and every old system solution is
obtained this way.

The literal sums of squares need a separate statement.
Every zero of the new sum lifts to a zero of the old sum
over every field, because the deleted row then vanishes and
the remaining rows agree.  The converse follows over fields
with a real embedding: a zero sum of squares forces each
row to vanish, and hence forces $z=X_{80}^2$.  Over a
general nonreal field, an old zero might instead arise by
cancellation among nonzero rows.  System elimination alone
does not prove that such a zero descends to the new sum.
We make no unrestricted converse assertion for literal sums.

Nevertheless the old generic $\mathbb Q(i)$ section has
$X_{80}=X_{81}=0$, so it does descend.  Its Gaussian
coordinate indices 96 and 97 become 95 and 96 after deletion.
Thus the new 1338-coordinate quartic, too, has input
projection equal to the whole field whenever $i$ belongs
to that field.  The script
\nolinkurl{check_quartic_base_change_20261106_v2.py}
verifies these new generic sections from the full sparse
tables and, independently, by squaring the restricted
residuals.  It writes a separate receipt, leaving the
preceding four certificates unchanged.

\subsection{Inlining three fixed constants}

The final ordinary-coefficient refinement substitutes the fixed
integers $3^{32}$, $3^{64}$, and $3^{128}$ for coordinates
58, 59, and 60 of the 1338-coordinate system, and deletes
their three defining rows.  The remaining system has 1335
coordinates and 1292 quadratic residuals.  These are private
polynomial graph equations, so substitution and reinsertion
give inverse maps on system solutions over every field.
As in the preceding elimination, the new literal sum of
squares always lifts to the old one; its converse follows
over fields with a real embedding, and is not asserted over
all nonreal fields merely from the system equivalence.

The largest newly introduced coefficient in a residual is
$3^{256}$, whose bit length is 406.  The associated new
quartic coefficient $3^{512}$ is less than $2^{812}$.
The full unchanged bound $H<2^{1020}$ therefore persists,
as also verified from the complete expanded coefficient
table.  The generic Gaussian section now uses coordinates
92 and 93, with all other witnesses zero.  Direct restriction
of the new table again gives $R(T)+a^2+b^2$ for an explicitly
recorded constant-plus-$T^4$ polynomial $R$; the same Gaussian
witness formulas prove the whole-field projection over any
characteristic-zero field containing $i$.  This section is
checked directly, rather than assuming that an old
cancellation solution survives fixed-constant substitution.

\subsection{The preferred degree-72 norm specifications}

All three transformed residual systems fit the same degree-36
cyclic division norm: their residual counts are 1272, 1295,
and 1292, respectively, each at most $36^2=1296$.
The resulting norm polynomials have 1315, 1338, and 1335 total
coordinates, respectively.  Their degree is exactly 72 and
their coefficient of $X_0^{72}$ is one, since the first
residual still contains $+X_0^2$ and all other residuals
are independent of the input coordinate.

The new translated residuals have only one coefficient
involving $S$, namely $8S$ at $U$ in the anchor row.  The
coefficient norm of every constant part is less than
$2^{511}<S$.  Hence every residual has coefficient norm
strictly less than
\[
                            9S=3^{2^{27}+2}.
\]
Since $3^5<2^8$,
\[
               \log_2(9S)<\frac85(2^{27}+2)=214748368.
\]
Combining this with $36!72^{36}<2^{361}$ gives the explicit
bound
\[
 H(P_{1315}^{\mathrm{norm,tr}})
    <36!72^{36}\,3^{36(2^{27}+2)}
    <2^{7730941609}.
\]
The ordinary 1338- and 1335-coordinate residuals still have coefficient
norm at most $2^{511}$, giving
\[
                       H(P_{1338}^{\mathrm{norm}}),\quad
                       H(P_{1335}^{\mathrm{norm}})
                                         \le2^{18757}.
\]
The complete specifications are
\nolinkurl{norm_repair_1315_translated_degree72_20261106_v2.json.gz}
\nolinkurl{norm_repair_1338_degree72_20261106_v2.json.gz},
and
\nolinkurl{norm_repair_1335_degree72_20261106_v2.json.gz}.
Each supplies the full residual list, cyclic matrix, and
cyclotomic coefficient-extraction rule.  Neither expands
the full degree-72 coefficient table or introduces a
quantified cyclotomic coordinate.

The fixed norm is faithful over $\mathbb Q(i)$ and over
extensions of degree coprime to 36, by the preceding local
and module arguments.  For any other specified number field,
the field-dependent prime construction supplies a faithful
degree-36 norm with the same coordinate counts and degree
72, with its stated prime-dependent coefficient bound.
Because these transformations preserve the residual
systems over all fields, these norm polynomials preserve
their input projections wherever the combining norms are
faithful.

In particular the unshifted scalar-contact contradictions
still hold for the transformed residual systems and their
faithful norm combinations.  Over a field with a real
embedding they also hold for the literal transformed sums
of squares.  These facts concern the unshifted artifacts;
they do not extend the contact counterexample to the
separately checked shift-2402 constructions.  Neither
transformation changes the unresolved arithmetic question
of defining $\mathcal O_K$ over a general number field.
